\chapter{Radiation: from optical properties to a thermal source}\label{radiation-theory}

Radiation changes the atmosphere by moving energy across its boundaries and
redistributing absorption and emission within it. It does not transport dry-air
mass. The selectable Nimbus families are \texttt{none}, \texttt{simple\_grey}
and \texttt{rrtmg}. The disabled family provides no radiative source. The other
two make substantially different optical approximations, even though both
ultimately provide temperature tendencies and flux diagnostics to the same
runtime owner.

\section{Define the sign, location and units before taking a divergence}

Let $F^\uparrow$ and $F^\downarrow$ be nonnegative upward and downward flux
magnitudes in $\mathrm{W\,m^{-2}}$, and define upward-positive net flux
$F=F^\uparrow-F^\downarrow$. Fluxes live at layer interfaces. In a flat column
with layer thickness $h_k$, the retained temperature conversion is
\[
 \dot T_k=-\frac{F_{k+1/2}-F_{k-1/2}}{\rho_{d,k}c_p h_k}.
\]
Here the runtime supplies dry density, and $c_p$ is the declared constant heat
capacity used by this radiation interface. An increase of upward net flux
through the layer means that energy leaves the layer, giving cooling. Reversing
the flux sign without reversing the divergence is a physical sign error, not
a choice of display convention.

Multiplying by dry layer mass per unit horizontal area gives
\[
 \rho_{d,k}c_p h_k\dot T_k=F_{k-1/2}-F_{k+1/2}.
\]
Summing over layers cancels every interior interface. The atmospheric power per
unit area is the bottom net flux minus the top net flux. This provides an
independent interface-flux check on the column tendency. It does not prove the
correctness of the optical coefficients or the chosen thermodynamic currency.
The \href{../src/radiation/column_evaluation.cpp}{column evaluator} checks this
closure before publishing the float theta tendency and cell power.

\begin{figure}[htbp]
\centering
\begin{tikzpicture}[x=1cm,y=1cm,font=\small,
  ray/.style={-{Latex[length=2mm]},very thick,draw=NimbusBlue}]
  \draw[fill=NimbusLight,draw=NimbusBlue] (0,0) rectangle (4,3);
  \draw[draw=NimbusBlue] (0,1)--(4,1) (0,2)--(4,2);
  \node[above] at (2,3.65) {Column interfaces carry flux};
  \draw[ray] (.8,3.5)--(.8,2.55);
  \node[left] at (.7,3.25) {$F^\downarrow$};
  \draw[ray] (3.2,2.55)--(3.2,3.5);
  \node[right] at (3.3,3.25) {$F^\uparrow$};
  \node at (2,2.5) {layer $k+1$};
  \node at (2,1.5) {layer $k$: $\rho_d,h,T$};
  \node at (2,.5) {layer $k-1$};
  \draw[very thick,draw=NimbusBlue] (-.2,0)--(4.2,0);
  \node[below=5pt] at (2,0) {surface: $T_s,\epsilon,\alpha$};
  \draw[ray] (4.15,1.5)--(5.15,1.5);
  \node[draw=NimbusBlue,fill=NimbusLight,rounded corners,align=center,
        text width=6cm,minimum height=2cm,anchor=west] at (5.3,1.5)
    {Interface difference\\[4pt]
     $\rho_d c_p h\dot T=F_{k-1/2}-F_{k+1/2}$\\[5pt]
     publish $\dot\theta$ and cell power together};
\end{tikzpicture}
\caption{Upward-positive net flux is $F=F^\uparrow-F^\downarrow$.
The layer gains energy when more net power enters its lower face than leaves
its upper face. Summing layers cancels interior exchanges. Optical closure,
surface boundary conditions and physical layer mass remain separate assumptions.}
\end{figure}


The conversion supplied to the dynamics owner is
$\dot\theta=\dot T(p_{\rm ref}/p)^{R_d/c_p}$ using pressure at the evaluated
state. The later EOS projection still participates in the thermal path described
in the coupling chapter. Column power, temperature tendency and the stored
theta increment are related quantities, not interchangeable numerical arrays.
For a terrain cell, physical heights, interface geometry, volume and surface
area must retain their declared meaning before applying an analogous budget.

\section{Optical depth summarizes an interaction along a path}

At a frequency $\nu$, let $I_\nu(\boldsymbol x,\boldsymbol\Omega)$ be radiance
in direction $\boldsymbol\Omega$. A local, stationary transfer equation with
elastic scattering can be written schematically as
\[
 \boldsymbol\Omega\cdot\nabla I_\nu
 =-\beta_{e,\nu}I_\nu+\beta_{a,\nu}B_\nu(T)
 +\frac{\beta_{s,\nu}}{4\pi}
    \int_{4\pi}P_\nu(\boldsymbol\Omega,\boldsymbol\Omega')
                      I_\nu(\boldsymbol\Omega')\,d\Omega'.
\]
Absorption and scattering coefficients have units $\mathrm{m^{-1}}$;
$\beta_e=\beta_a+\beta_s$. This convention normalizes the phase function so
its angular integral is $4\pi$. The thermal term assumes local thermodynamic
equilibrium for the emitting material. A model that drops the scattering term
has made a different physical approximation from one that replaces its angular
integral with a quadrature.

For pure absorption along path coordinate $s$, optical depth is
$\tau=\int\beta_a\,ds$. With no emission on that path, the solution is
$I=I_0e^{-\tau}$. A vertical optical depth $\Delta\tau$ seen by a direct beam
with vertical cosine $\mu_0>0$ becomes path depth $\Delta\tau/\mu_0$ in a
plane-parallel layer. This geometric factor belongs to the beam direction.
A diffuse hemisphere contains many path angles and cannot generally use the
same single cosine.

Nimbus's schemes reduce spectral and angular detail in different ways. The
simple-grey scheme uses broadband optical depths. RRTMG uses spectral bands and
reduced absorption-distribution quadrature. The underlying transfer equations
help identify assumptions, but the code's quadrature, boundary rules and
coefficient inputs define the numerical calculation actually performed.

\section{The two simple-grey optical models}

The pressure-only clear-sky model defines an interface optical-depth profile
\[
 \tau(p)=\tau_{\rm ref}
       \left(\frac{p}{100000\,\mathrm{Pa}}\right)^n,
 \qquad\Delta\tau_k=\tau(p_{k-1/2})-\tau(p_{k+1/2}).
\]
Longwave and shortwave have separate reference depths and exponents. Pressure
interfaces must be finite, nonnegative and nonincreasing upward, with positive
surface pressure. This is a prescribed broadband profile, not a calculation
of molecular absorption from a complete gas spectrum. It deliberately provides
a simple radiative experiment whose parameters are model assumptions.

The absorbing water/cloud-path model adds absorber mass paths,
\[
 \Delta\tau_k^{\rm water}
 =\rho_{d,k}h_k
   \big(\kappa_vq_{v,k}+\kappa_lq_{c,k}+\kappa_iq_{i,k}\big).
\]
The mass-absorption coefficients $\kappa_x$ have units
$\mathrm{m^2\,kg^{-1}}$, so the expression is dimensionless. Separate
coefficients are supplied for the two broadband components. This path uses
vapor, cloud liquid and cloud ice. It does not automatically turn rain, snow,
graupel or hail into optical contributors merely because those categories
exist in microphysics.

The absorbing model adds this term to the pressure-profile layer depth. It is
not a scattering cloud model: changing cloud water changes absorption, but
does not create the cloud scattering and angular redistribution represented by
a more elaborate optical scheme. Configuration must pair the selected optical
model with its clear-sky setting consistently. The
\href{../src/radiation/schemes/simple_grey/simple_grey.cpp}{simple-grey owner}
validates those choices and requires physical column inputs.

\section{Grey longwave transfer and the surface boundary}

The implemented longwave calculation treats each layer as isothermal for its
broadband source. With $\mathcal T_k=e^{-\Delta\tau_k}$, it uses recurrences
of the form
\[
 \begin{aligned}
 F^\downarrow_{k-1/2}
 &=\mathcal T_k F^\downarrow_{k+1/2}
   +(1-\mathcal T_k)\sigma T_k^4,\\
 F^\uparrow_{k+1/2}
 &=\mathcal T_k F^\uparrow_{k-1/2}
   +(1-\mathcal T_k)\sigma T_k^4.
 \end{aligned}
\]
These equations describe the code's effective optical-depth convention; no
additional longwave diffusivity factor should be inserted when reproducing this
particular grey calculation. They assume a broadband thermal source and neglect
longwave scattering in this path.

The top incoming longwave flux is zero. After the downward sweep reaches the
surface, the upward lower boundary is
\[
 F^\uparrow_s=\epsilon\sigma T_s^4
                 +(1-\epsilon)F^\downarrow_s.
\]
The second term represents reflected downwelling longwave. Omitting it would
make a nonblack surface absorb energy that its declared emissivity does not
permit under this opaque-surface convention. The upward sweep therefore depends
on the completed downward sweep. Surface temperature and emissivity are explicit
inputs; the radiation calculation does not solve a soil or land-surface energy
balance that prognoses a new $T_s$.

The recurrences and the flux-divergence conversion are separate stages in
\href{../src/radiation/base/radiative_transfer.cpp}{radiative transfer}.
A valid recurrence can still produce the wrong heating if layer indexing,
interface orientation or the mass used by the conversion is inconsistent.

\section{Direct sunlight and Lambertian reflection are different angular fields}

For an illuminated flat column, the top incident shortwave flux is
$F^\downarrow_{\rm top}=\mu_0 S_0$, where $S_0$ is irradiance normal to the
solar beam. The downward grey path attenuates each layer by
$e^{-\Delta\tau_k/\mu_0}$. The nighttime branch publishes no shortwave source
rather than evaluating a singular or negative path length at $\mu_0\le0$.

At the surface, the reflected upward flux is
$F^\uparrow_s=\alpha F^\downarrow_s$. Nimbus treats this reflected boundary
as Lambertian. Its radiance is direction-independent over the upward hemisphere,
so the upward flux at cumulative optical depth $\tau$ above the surface is
\[
 F^\uparrow(\tau)=F^\uparrow_s\,2E_3(\tau),\qquad
 E_3(\tau)=\int_1^\infty\frac{e^{-\tau t}}{t^3}\,dt.
\]
The factor follows directly from hemispheric integration:
$2\int_0^1\mu e^{-\tau/\mu}\,d\mu=2E_3(\tau)$.
It approaches one at zero optical depth. Applying the direct solar cosine to
this reflected flux would give it the wrong angular distribution.

This grey construction contains atmospheric absorption and a reflecting lower
boundary but no atmospheric scattering. Albedo changes the energy partition
between the surface and the upward radiation; it is not a source of energy.
The implementation evaluates the exponential-integral transmission with bounded
numerical branches and rejects nonphysical results. Tests of the thin, thick,
black-surface and reflecting-surface limits examine different parts of this
contract.

\section{What the correlated-k approximation changes}

Molecular absorption varies rapidly with frequency. Resolving every spectral
line in every grid column would be expensive. Within a band, an absorption
coefficient distribution can be represented by a cumulative variable $g$,
turning a band transmission integral into an approximation such as
\[
 \overline{\mathcal T}(u)
 =\int_0^1e^{-k(g)u}\,dg
 \simeq\sum_j w_j e^{-k_j u}.
\]
Here $u$ is the absorber amount in the coefficient convention and the weights
cover the distribution interval. Using correlated ordering between layers,
combining gases and interpolating coefficients across temperature and pressure
introduce additional modeling assumptions. This is a spectral quadrature;
its $g$ coordinate is not gravitational acceleration or a particle-size moment.

\href{https://doi.org/10.1029/97JD00237}{Mlawer et al. (1997)} develops the
longwave correlated-k reference method, and
\href{https://doi.org/10.1029/2008JD009944}{Iacono et al. (2008)} describes the
AER model family and its greenhouse-gas forcing calculations. Their validation
results do not automatically apply to Nimbus's translated implementation,
coefficient versions, cloud assumptions or coupled runtime.

Nimbus implements the calculation in C++; the AER netCDF files provide
coefficients rather than an executable Fortran solver. Typed loading validates
the expected coefficient schema and identity. Pressure/temperature preparation,
major and minor gas absorption, continua, cloud optics and transfer have
separate source owners under
\href{../src/radiation/schemes/rrtmg/}{the RRTMG implementation}. The reduced
spectral structure is part of those coefficient contracts, not a tunable number
of arbitrary equal-frequency samples.

\section{Gas columns, cloud optical properties and transfer closure}

Nimbus constructs dry-air molecular columns from the supplied dry density and
physical layer thickness. With dry molar mass $M_d$ in $\mathrm{kg\,mol^{-1}}$,
\[
 N_d=\frac{\rho_d h}{M_d}\,\frac{N_A}{10^4}
\]
gives molecules per square centimetre. A water dry-air mass mixing ratio is
converted to a dry-air-relative molar ratio with
$r_v=q_v M_d/M_v$. Gas mole fractions and their allowed totals have explicit
validation. Pressure, temperature and molecular column amounts therefore have
different roles; substituting a pressure-thickness formula for the actual
$\rho_d h$ preparation without reviewing its assumptions changes the inputs.

The runtime currently passes cloud liquid $q_c$ and cloud ice $q_i$ to the
optics. In cloudy mode it uses liquid/ice mass paths
$1000\rho_d h q_c$ and $1000\rho_d h q_i$ in $\mathrm{g\,m^{-2}}$, with
configured effective radii in micrometres. A cell with either positive cloud
category is assigned cloud fraction one; otherwise it is clear. There is no
prognostic subgrid cloud fraction or general cloud-overlap model in this
construction. Effective radii are optical inputs here, not automatically the
radius diagnosed from every microphysics moment. Falling condensate categories
are not silently summed into the two cloud paths.

Longwave transfer includes its band-dependent diffusivity treatment. Shortwave
transfer uses the implemented practical improved flux method (PIFM) two-stream
layer coefficients, direct/diffuse reflection and transmission, and the
supplied scattering/asymmetry information. A two-stream solution represents an
angular closure, not the full radiance field at every direction. Conservative
scattering and optically thin limits need stable numerical branches because
algebraically equivalent expressions can suffer cancellation or singular-looking
intermediate terms. See
\href{../src/radiation/schemes/rrtmg/rrtmg_longwave_transfer.cpp}{longwave transfer}
and \href{../src/radiation/schemes/rrtmg/rrtmg_shortwave_transfer.cpp}{shortwave transfer}.

Interface temperatures, when not supplied, use the retained pressure-coordinate
interpolation and endpoint rules. The physical column order and the transfer
solver's sweep order are explicitly converted. Agreement of a single surface
flux does not establish that all interface fluxes, layer heating and cloud/gas
responses agree with an independent column reference.

\section{Terrain adds visibility, projected area and horizontal exchange}

For a unit direction toward the Sun, $\boldsymbol s$, and an upward surface
normal $\boldsymbol n$, direct intercepted irradiance is proportional to
$\max(0,\boldsymbol n\cdot\boldsymbol s)$. Visibility is an additional condition:
a face aimed toward the Sun may still be shadowed by intervening terrain.
Area measured along the slope differs from horizontal plan area. Fluxes in
$\mathrm{W\,m^{-2}}$ must be paired with the area their definition uses before
becoming power in watts.

Nimbus has explicit terrain shortwave paths for visibility, ray traversal,
absorbed power, reflected angular transport and spectral variants. A ray through
several physical cells carries an energy amount; attenuation along each segment
partitions it between absorption and outgoing power. A reflecting terrain
boundary can redirect it into other columns. Reusing independent vertical-column
heating while merely shading a picture would not represent that exchange.

For a bounded shortwave transport calculation, an appropriate whole-domain
partition is
\[
 P_{\rm in}=P_{\rm atmosphere}+P_{\rm surface}
        +P_{\rm top\ escape}+P_{\rm lateral\ escape}+P_{\rm unresolved}.
\]
The unresolved term records power remaining when a declared angular/bounce or
convergence bound is reached; it is not automatically atmospheric heating.
Changing quadrature, reflection count or stopping tolerance changes a numerical
approximation that must be checked independently of this bookkeeping identity.

Source presence does not imply every terrain/backend combination is admitted.
The current runtime rejects mapped RRTMG shortwave on an active accelerator
because its native spectral-terrain owner is missing. CPU mapped spectral work
and native column batches have separate scopes. The
\href{../src/radiation/runtime.cpp}{radiation runtime} is the authority for the
actual selected path and required rejection behavior.

\section{Cadence, stored power and the complete coupled step}

Radiation is refreshed on its declared cadence and can hold its source between
calls. The evaluated optical state, published heating and held power must share
one completed generation. The next dynamics interval should not combine a newly
computed power diagnostic with an older held theta tendency. A refresh failure
must preserve the preceding coherent runtime state and its timing.

If a held atmospheric power is $P_h$ over an interval $\Delta t$, its declared
energy input is $P_h\Delta t$. The
\href{../src/radiation/energy_ledger.cpp}{radiation energy ledger} accounts for
that held input and the relevant source endpoint. This is distinct from simply
recomputing an interface-flux difference from a later atmospheric column.
Cadence holding approximates time evolution; a smaller dynamics step does not
by itself make a long-held radiation source current.

The checks in \href{../tests/physics/test_radiation.cpp}{radiation tests} and
\href{../tests/physics/test_radiation_runtime_state.cpp}{runtime-state tests}
cover columns, units, invalid inputs, closure and publication behavior for their
stated scopes. Broader coupled, terrain, restart and native paths need their own
bounded and campaign evidence. Neither a familiar scheme name nor a finite
flux residual replaces that validation.
